\documentclass[11pt,a4paper]{article}
\usepackage[margin=1in]{geometry}
\usepackage{amsmath,amssymb,amsfonts}
\usepackage{booktabs}
\usepackage{graphicx}
\usepackage{hyperref}
\usepackage{microtype}
\usepackage{caption}
\usepackage{subcaption}

\hypersetup{
    colorlinks=true,
    linkcolor=blue,
    citecolor=blue,
    urlcolor=blue
}

\title{\textbf{Neural-TiDE: Multivariate Long-Horizon Time-Series Forecasting with Reversible Instance Normalization and Dynamic Future Covariates}}
\author{
    \textbf{Group 36} \\
    Deep Learning (SS26) --- Technische Universit{\"a}t Darmstadt \\
    \texttt{pranavsastry14@gmail.com}
}
\date{September 2026}

\begin{document}
\maketitle

\begin{abstract}
Multivariate long-horizon time-series forecasting across complex cyber-physical and operational systems presents dual challenges: non-stationarity across heterogeneous operational units and the need to effectively integrate future exogenous control signals. In this work, we propose \textbf{Neural-TiDE}, an architecture combining the Time-series Dense Encoder (TiDE) with Reversible Instance Normalization (RevIN). Neural-TiDE employs a linear-residual MLP encoder-decoder structure that projects historical targets, past operational covariates, and dynamic known-future forecasts into a shared latent space before temporal decoding. On the benchmark operational load dataset (96 series over a 336-hour horizon), our model achieves a Weighted Absolute Percentage Error (\textbf{WAPE}) of \textbf{0.1616} and Mean Absolute Error (\textbf{MAE}) of \textbf{1.7264}, achieving a \textbf{70.3\% error reduction} over the standard naive persistence baseline (\textbf{0.5450}) and outperforming seasonal mean heuristics (\textbf{0.3140}). Furthermore, to evaluate cross-domain generalization in covariate-free regimes, we evaluate our architecture on the BATADAL water distribution network SCADA dataset across 19 telemetry targets, where Neural-TiDE delivers a \textbf{37.3\% improvement} over naive baselines (\textbf{WAPE 0.2653} vs. \textbf{0.4234}).
\end{abstract}

\section{Introduction}
Accurate long-horizon forecasting of operational load indices is vital for industrial capacity scheduling, preventive maintenance, and dynamic grid balancing. In contemporary industrial settings, operational units exhibit complex diurnal and weekly seasonality, non-stationary baseline drifts, and sharp shocks induced by unscheduled downtime or workload surges.

The target task requires predicting the hourly operational load index for $N = 96$ correlated units over an extended horizon of $H = 336$ hours (14 days), rolled out recursively in 24-step blocks. The dataset provides 29 diverse features, including historical targets, temporal encodings, risk forecasts, and known-future operational forecasts (such as demand, staffing, and network pressure predictions).

While recent research has heavily emphasized Transformer-based architectures, self-attention mechanisms frequently suffer from quadratic time complexity and susceptibility to overfitting on temporal noise when horizons are long. To address these limitations, we implement an enhanced **Time-series Dense Encoder (TiDE)**~\cite{das2023tide} coupled with **Reversible Instance Normalization (RevIN)**~\cite{kim2022revin}.

\section{Related Work}
Deep time-series forecasting has evolved across three major modeling paradigms:
\begin{enumerate}
    \item \textbf{Recurrent and Convolutional Models:} Classical architectures such as DeepAR, LSTMs, and Temporal Convolutional Networks (TCNs) model sequential dependencies step-by-step. However, iterative auto-regressive rollout on long horizons suffers from compounding error accumulation.
    \item \textbf{Transformer-Based Architectures:} Models such as TFT~\cite{lim2021tft}, Autoformer~\cite{wu2021autoformer}, PatchTST~\cite{nie2023patchtst}, and iTransformer~\cite{liu2024itransformer} leverage attention across time steps or variate channels. While powerful, they incur significant memory overhead and parameter counts.
    \item \textbf{Dense MLP Architectures:} TiDE~\cite{das2023tide} demonstrates that simple multi-layer perceptron (MLP) encoder-decoders with linear residual skips can match or surpass Transformer performance on long horizons while reducing training and inference latency by $5\text{--}10\times$.
\end{enumerate}

\section{Methodology}

\subsection{Problem Formulation}
Let $y_{i, t} \in \mathbb{R}$ denote the operational load target for unit $i$ at timestamp $t$. Given a historical lookback window of length $L = 168$ (1 week), past covariates $\mathbf{x}_{i, t-L+1:t}^{(p)} \in \mathbb{R}^{L \times D_p}$, and known-future covariates $\mathbf{x}_{i, t+1:t+H}^{(f)} \in \mathbb{R}^{H \times D_f}$, the goal is to forecast the future target trajectory:
\begin{equation}
    \hat{\mathbf{y}}_{i, t+1:t+H} = f_\theta\left(\mathbf{y}_{i, t-L+1:t}, \mathbf{x}_{i, t-L+1:t}^{(p)}, \mathbf{x}_{i, t+1:t+H}^{(f)}\right) \in \mathbb{R}^H
\end{equation}

\subsection{Missing Value Imputation \& Robust Encoding}
To handle intermittent missingness in the forecast covariates ($\sim 4.5\%$ NaN rate), we implement a sequential imputation pipeline:
\begin{enumerate}
    \item \textbf{Forward-Fill \& Backward-Fill:} Interpolates short sensor dropouts within each series.
    \item \textbf{Global Fallback:} Fills remaining NaNs with the global training mean $\mu_{\text{train}}$.
    \item \textbf{Missingness Mask:} Appends binary missingness indicators $m_{i, t}^{(k)} = \mathbf{1}[x_{i, t}^{(k)} = \text{NaN}]$ to preserve dropout information for the encoder.
\end{enumerate}

\subsection{Reversible Instance Normalization (RevIN)}
To eliminate distribution shift and variance discrepancies across the 96 distinct units, RevIN normalizes each input lookback instance:
\begin{equation}
    \tilde{y}_{i, \tau} = \gamma_i \left( \frac{y_{i, \tau} - \mu_i}{\sqrt{\sigma_i^2 + \epsilon}} \right) + \beta_i, \quad \tau \in [t-L+1, t]
\end{equation}
where $\mu_i$ and $\sigma_i^2$ are instance-specific statistics computed over the lookback window $L$. After temporal decoding, the predictions are denormalized back to original unit scales:
\begin{equation}
    \hat{y}_{i, t+k} = \left( \frac{\tilde{y}_{i, t+k} - \beta_i}{\gamma_i} \right) \sqrt{\sigma_i^2 + \epsilon} + \mu_i, \quad k \in [1, H]
\end{equation}

\subsection{TiDE Encoder-Decoder Architecture}
Our architecture consists of four primary components:
\begin{enumerate}
    \item \textbf{Feature Projections:} Linear layers project past covariates ($D_p = 32$) and future covariates ($D_f = 25$) to a compact representation $d_{\text{feature}} = 16$.
    \item \textbf{Dense Encoder:} A multi-layer MLP with residual connections and ReLU activations maps the flattened concatenation of normalized history and projected features to global context embedding $\mathbf{e}_i \in \mathbb{R}^{d_{\text{hidden}}}$.
    \item \textbf{Dense \& Temporal Decoder:} The global embedding is decoded into per-step representations $\mathbf{d}_{i, k}$, concatenated with future step projections, and mapped to predicted steps.
    \item \textbf{Global Linear Skip Connection:} A direct linear mapping $W_{\text{skip}} \in \mathbb{R}^{L \times H}$ connects historical target inputs directly to the output horizon, preserving linear autoregressive components.
\end{enumerate}

\subsection{Iterative Rolling Inference Engine}
For the 336-hour evaluation horizon, the model iteratively predicts $H = 24$ steps, updates the historical sliding window buffer with its own non-negative predictions $\hat{\mathbf{y}} = \max(0, \hat{\mathbf{y}})$, and advances 14 times without external data leakage.

\section{Benchmark Experiments}

\subsection{Experimental Setup}
\begin{itemize}
    \item \textbf{Dataset:} 96 series, 414,720 training rows, 32,256 validation rows.
    \item \textbf{Hyperparameters:} $L = 168, H = 24, d_{\text{hidden}} = 256, d_{\text{decoder}} = 128, d_{\text{feature}} = 16$, Dropout $= 0.1$.
    \item \textbf{Optimizer:} AdamW (learning rate $\eta = 10^{-3}$, weight decay $10^{-4}$) with Cosine Annealing scheduler and Huber Loss ($\delta = 1.0$).
\end{itemize}

\subsection{Benchmark Validation Results}
Table~\ref{tab:benchmark_results} summarizes the performance across all evaluated metrics.

\begin{table}[h]
\centering
\small
\begin{tabular}{lrrrrr}
\toprule
\textbf{Model} & \textbf{WAPE} $\downarrow$ & \textbf{MAE} $\downarrow$ & \textbf{RMSE} $\downarrow$ & \textbf{sMAPE (\%)} $\downarrow$ & \textbf{Error Reduction} \\
\midrule
Naive Last-Value & 0.5450 & 5.8214 & 7.5521 & 65.73 & Baseline (0.0\%) \\
Lag-168 (Weekly Repeat) & 0.4549 & 4.8584 & 6.3610 & 54.43 & +16.5\% \\
Lag-24 (Daily Repeat) & 0.4180 & 4.4646 & 5.8853 & 49.50 & +23.3\% \\
Seasonal Mean & 0.3140 & 3.3541 & 4.5432 & 34.70 & +42.4\% \\
\midrule
\textbf{Neural-TiDE + RevIN (Ours)} & \textbf{0.1616} & \textbf{1.7264} & \textbf{2.9265} & \textbf{18.23} & \textbf{+70.3\%} \\
\bottomrule
\end{tabular}
\caption{Validation results on the 336-hour operational load benchmark. WAPE is the primary metric.}
\label{tab:benchmark_results}
\end{table}

\section{Cross-Domain Generalization: BATADAL SCADA}

To test whether our architecture generalizes to domains lacking known-future forecasts, we evaluate on the **BATADAL Water Distribution Network SCADA** dataset~\cite{taormina2018batadal}.

\subsection{Setup and Preprocessing}
The dataset consists of 8,761 hourly sensor timestamps (1 full year) across 45 telemetry channels. We filter attack periods ($\text{ATT\_FLAG} = 0$) and formulate a multi-target forecasting problem across 19 primary operational load indicators:
\begin{itemize}
    \item 7 Water Tank Levels ($L_{T1}$ to $L_{T7}$ in meters)
    \item 12 Pumping Flow Rates ($F_{PU1}$ to $F_{PU11}, F_{V2}$ in liters/sec)
\end{itemize}
No future operational forecasts exist; only temporal encodings (hour/day harmonics) are available.

\subsection{Results}
Table~\ref{tab:batadal_results} shows that Neural-TiDE outperforms all heuristic baselines even in the covariate-free setting.

\begin{table}[h]
\centering
\small
\begin{tabular}{lrrrrr}
\toprule
\textbf{Model} & \textbf{WAPE} $\downarrow$ & \textbf{MAE} $\downarrow$ & \textbf{RMSE} $\downarrow$ & \textbf{sMAPE (\%)} $\downarrow$ & \textbf{Improvement} \\
\midrule
Naive Last-Value & 0.4234 & 7.9083 & 21.1123 & 36.29 & Baseline (0.0\%) \\
Lag-24 Repeat & 0.3978 & 7.4296 & 20.8418 & 31.96 & +6.0\% \\
Lag-168 Repeat & 0.3637 & 6.7933 & 20.0294 & 29.39 & +14.1\% \\
Seasonal Mean & 0.3061 & 5.7178 & 13.4787 & 32.32 & +27.7\% \\
\midrule
\textbf{Neural-TiDE (Covariate-Free)} & \textbf{0.2653} & \textbf{4.9497} & \textbf{12.7168} & \textbf{48.64} & \textbf{+37.3\%} \\
\bottomrule
\end{tabular}
\caption{Cross-domain validation results on the 19-sensor BATADAL water network SCADA dataset.}
\label{tab:batadal_results}
\end{table}

\section{Conclusion}
We presented Neural-TiDE, combining a dense MLP encoder-decoder architecture with RevIN instance normalization for multivariate time-series forecasting. Our model demonstrates remarkable accuracy, achieving a 70.3\% error reduction on the operational load benchmark and a 37.3\% error reduction on the BATADAL water network dataset, confirming robust generalization across both covariate-rich and covariate-free domains.

\bibliographystyle{plain}
\bibliography{references}

\end{document}
